% Generated from validated synthesis. Do not hand-edit.
\section*{Monthly Signals}
\addcontentsline{toc}{section}{Monthly Signals}
\smallskip
\noindent\textbf{Agentはmodelから実行系へ} harness、context、desktop、terminal、computer useが別々の実装面で前景化し、Agentの品質をbase modelだけでは説明できなくなり始めた。 \hfill{\footnotesize p.~\pageref{pkg:agent-stack-leaves-chat}}\par
\smallskip
\noindent\textbf{外部actionには制御境界が要る} webやtoolへの接続がprompt injectionをdata flowと外部actionの問題へ変え、automatic executionとuser confirmationの境界がarchitecture上の課題になった。 \hfill{\footnotesize p.~\pageref{pkg:autonomy-needs-guardrails}}\par
\smallskip
\noindent\textbf{訓練環境を合成する} world modelとsynthetic toolsという異なる方法で、tool-using Agentのtrainingをlive environmentから切り離し、cost・privacy・stabilityを扱う方向が示された。 \hfill{\footnotesize p.~\pageref{pkg:training-agents-offline}}\par
\smallskip
\noindent\textbf{Tool-native化はmodelからruntimeまで続く} retrievalやcode executionを前提とするmodel/APIと、新しいmodel familyやworkloadを短い周期で取り込むserving runtimeが、同じstackの上下で進んだ。 \hfill{\footnotesize p.~\pageref{pkg:tool-native-models} / p.~\pageref{pkg:runtime-catches-the-agent-wave}}\par
\smallskip
\noindent\textbf{後続月を読むための出発点} 1月はagentic engineeringの完成ではなく、harness、security、training、tool-enabled API、runtimeという観察軸が立ち上がった月として位置付ける。 \hfill{\footnotesize p.~\pageref{pkg:closing-synthesis}}\par
\medskip
\begin{claimboundary}[Retrospective scope]
本号は2026年1月を後日確認可能になった一次情報も用いて再構成するRetrospective Specialである。Coverage Periodと制作時点を同一視せず、vendor / project / author claimの境界は各記事とTechnical Notesで明示する。
\end{claimboundary}
\medskip
\setcounter{tocdepth}{1}
\tableofcontents
